%% PASS 3: Chapter 1 compared with Dablander's DPhil thesis, Chapter 1 (Introduction, pp. 1-6, with section 1.1 "Outline of Thesis and Research Contributions").
%% His introduction is three pages of prose and a three-page chapter-by-chapter outline. It contains no percentages, counts or parameter lists.
%% His outline gives each chapter what it does, its findings in words, one sentence on what is new, and where it is published.
%% He has no separate research-questions or collaborations section; funding and Lhasa sit in his Acknowledgements. Her sections stay, since they are her headings.
%% Cut here: the numbers planned for Section 1.4, the lists of distances and representations for the opener, and the named comparison methods for Chapter 5.
%% Added: one bullet asking for a "what is new" sentence per chapter, which his outline has and these notes lacked.

% ==== FROM: Introduction notes, Paragraph 4 ====
% - Change: "vector operations" becomes the three distances tried: Tanimoto on ECFP4, Euclidean on PDV, and Mahalanobis.
% - Change: "representations" becomes the five tried: ECFP4, Sort \& Slice, PDV, and canonical and randomised SMILES.
%   (Replaced by: "vector operations" becomes "distances"; the lists go in Section~\ref{sec:lnl_motivation}.)

% ==== FROM: From noise robustness to adaptive representations, Paragraph 1 ====
% - Under Gaussian noise on QM9, model architecture explains 49.5% of the variance in robustness. Representation explains 9.2%.
% - Under censoring, models were most confident on the clipped labels.

% ==== FROM: From noise robustness to adaptive representations, Paragraph 2 ====
% - ... tested in 221 configurations ...
% - All 20 configurations with the lowest retention of R^2 judged a label by its neighbours' labels.

% ==== FROM: From noise robustness to adaptive representations, Paragraph 3 ====
% - The growth is linear in the number of features, not exponential.

% ==== FROM: Outline of the thesis and contributions, Chapter 5 ====
% - Will test why: against PCA, mRMR, recursive feature elimination and random features at the same number of features.
